\section{Related Work}
\label{sec:related}

\nospacestitle{LLM requests global scheduling. \,}
%Efficiently scheduling requests in cluster is non-trivial due to the variance of the upcoming workloads and the complex trade-offs between different SLO metrics.
{\sys} continues the line of research on scheduling LLM requests in a cluster~\cite{vllm,DBLP:journals/corr/abs-2506-05871,DBLP:conf/iclr/SrivatsaHAL025,
305212,DBLP:journals/corr/abs-2507-17769,llm-d}.
To the best of our knowledge, all these methods attempt
to achieve both {\kvcache}-awareness and load-balancing
through three approaches in combining different indicators for each objective,
as we have extensively discussed in \textsection{\ref{sec:principles}}.
{\sys} builds on these works to reuse their indicators but further
proposes a new and simple multiplication combinator,
and our extensive evaluations confirm the effectiveness of our approach.

\iffalse
Existing works can be broadly into three groups: 
1. Linear-Combination-based: vLLM~\cite{vllm} and BestServe~\cite{DBLP:journals/corr/abs-2506-05871} 
both aim to balance running requests, 
but vLLM additionally accounts for waiting requests, 
while BestServe choose cumulative requests
2. Filter-based: Preble~\cite{DBLP:conf/iclr/SrivatsaHAL025} filtered out 
instance when batch-size range exceeds the specified threshold, 
and then select instance with the highest kvcache hit.
3. SLO-based scheduling: Mooncake \cite{305212} predict the TTFT 
for each prefill instances and take KVCache transfer cost into 
consideration. PolyServe~\cite{DBLP:journals/corr/abs-2507-17769} estimated TPOT 
and schedule requests to instance of the highest TPOT without 
SLO-violation for elastically scaling. While our work used a simple
multiplication of system metrics, avoiding complex modeling, 
combining and filtering, capturing the essential factors affecting
SLOs in colocation scenarios and worked well.
\fi

\stitle{LLM requests scheduling within an instance. \,} 
Besides global scheduling across instances,
a line of work focuses on local scheduling optimizations within an instance.
Sarathi-Serve~\cite{298679} introduces chunked prefill, which splits long prefill requests into smaller chunks
to reduce stall time for co-located decode requests.
VTC~\cite{DBLP:conf/osdi/0007CLZ0ZGS24} adopts a token-based admission control mechanism
to achieve fairness.
FairBatching~\cite{DBLP:journals/corr/abs-2510-14392} uses a linear-time analytical model
to prioritize prefill versus decode tokens and dynamically form batches.
These works are orthogonal to our global scheduling, and
a good global scheduler can further enhance local scheduling effectiveness,
e.g., by reducing the frequency of overloading that is non-trivial to handle with local scheduling alone.

\stitle{Optimizing LLM serving. \,}
Besides scheduling,
a variety of techniques have been proposed to optimize LLM serving performance.
These include methods
to improve {\kvcache} hit rates~\cite{305212,10.5555/3768039.3768067},
to provide elasticity for model serving~\cite{DBLP:conf/osdi/ZhangWLWS0025,DBLP:conf/sosp/Xiang0QYZYZL0025},
and to optimize GPU execution efficiency~\cite{vllm,DBLP:conf/asplos/KamathPM0RP25},
to name a few.
To the best of our knowledge, all these techniques coexist with scheduling optimizations,
so they can be combined with {\sys} to further improve overall performance.
\iffalse
%while Mooncake~\cite{305212} further implements an RDMA-based 
In serverless deployments, BlitzScale~\cite{DBLP:conf/osdi/ZhangWLWS0025} developed rapid autoscaling mechanisms 
to elastically provision LLM instances. 
PagedAttention~\cite{vllm} focused on memory management by organizing KV cache in pages, 
enabling higher concurrency under limited GPU memory, 
whereas POD-Attention~\cite{DBLP:conf/asplos/KamathPM0RP25} improved chunked prefill implementations 
at the kernel level by sharing GPU compute resources between 
prefill and decode operations. 
These techniques are orthogonal to our scheduling focus and 
could be combined with our global scheduler for additional performance gains.
\fi